\section{I/O performance}
\label{section:high-level:io}

\begin{summary}
 We consider a code not to suffer from IO constraints if it spends less than 20\% of its runtime in I/O.
\end{summary}

\subsection{Benchmark preparation}

%When preparing the benchmark for core, intra-node, inter-node and GPU analysis, we advised turning off I/O where possible as this can make some of the measurements unreliable. However, I/O performance is crucial for some codes and in some cases necessary for the code to run. 

When benchmarking I/O, or benchmarking with I/O enabled, we advise using a benchmark with a typical I/O load - this should be verified by the code's developers. It can tempting to reduce the I/O operations to `game' this metric, but if a code does rely on heavy and frequent I/O then this should be included in the characteristic benchmark.

\subsection{Data Collection}

There are relatively few tools that focus on I/O performance. Hence, we are restricted in the number of tools that can collect our metric of interest, namely, the proportion of runtime spent doing I/O operations. We list below an example of how to do this with a Darshan, an I/O focused performance tool.

\begin{instructions}{Darshan}
In this example we present paths and modules as they are defined on a local system. Any user of this book would therefore have to find the relevant Darshan files on their system of choice.

For a dynamic runtime analysis we setup Darshan with
\begin{verbatim}
export DARSHAN_DIR=/apps/developers/tools/darshan/3.3.1/1/

intel-2021.4-intelmpi-2021.6 # path to Darshan
export DARSHAN_LOGDIR=./darshan_logs
mpiexec -genv LD_PRELOAD ${DARSHAN_DIR}/lib/libdarshan.so
-np 4 ./code
\end{verbatim}
in which the \texttt{darshan\_logs} directory is user defined for hosting Darshan profile outputs. The logging environment variable name can be different so make sure to run
\begin{verbatim}
    darshan-config --log-path
\end{verbatim}
to find the name on your system.
Once this has created a \texttt{*.darshan} binary file this can analysed via
\begin{verbatim}
    darshan-parser --total my_log.darshan 
\end{verbatim}

The relevant metrics from this output are in the form
\begin{verbatim}
STDIO_F_*_TIME: cumulative time spent in different 
types of functions.
\end{verbatim}
Therefore depending on different forms of I/O --- e.g., \texttt{POSIX}, \texttt{MPIIO}, \texttt{STDIO}, \texttt{H5F}, \texttt{H5D}, \texttt{PNETCDF\_FILE}, \texttt{LUSTRE}, \texttt{DFS}, \texttt{DAOS} --- we need to sum together the read and write times.
\end{instructions}

%\begin{instructions}{VTune}
% t.b.d.
%\end{instructions}

\subsection{Evaluation}

\begin{performancemetric}
 We evaluate 
 \[
  C_{\text{I/O}} = 1 - \frac{t_{\text{read}} + t_{\text{write}}}{t_{\text{total}}}
\]
where $t_{\text{read}}$ and $t_{\text{write}}$ are the times the code spends in reading and writing data, respectively, whilst $t_{\text{total}}$ is simply the total runtime (inclusive of $t_{\text{read}}$ and $t_{\text{write}}$).
\end{performancemetric}

\begin{center}
 \begin{tabular}{p{0.34\textwidth}p{0.62\textwidth}}
  I/O Proportion & Description \\
  \hline
  $C_{\text{I/O}} \geq 0.8$
   & \cellcolor{green!25} I/O does not dominate runtime.
  \\
  $0.6 \leq C_{\text{I/O}} < 0.8$
   & \cellcolor{yellow!25} I/O forms a significant part of the total runtime.
  \\
  $C_{\text{I/O}} < 0.6$
   & \cellcolor{red!25} I/O dominates the runtime.
 \end{tabular}
\end{center}


\paragraph*{How it works}

\begin{itemize}[noitemsep]
  \item[$\boxplus$] Run an indicative I/O workload and observe time spent in reads and writes
  \item[$\boxplus$] Accumulate time from different forms of reads and writes into a single I/O runtime
  \item[$\boxplus$] Observe how much of total runtime is spent in this I/O runtime
\end{itemize}


\paragraph*{How it does not work}

\begin{itemize}[noitemsep]
  \item[$\boxminus$] Run your benchmark without any indicative I/O workload
  \item[$\boxminus$] Dig into the details of the code's I/O from tools like Darshan, and begin to draw conclusions
\end{itemize}
